% GENERATED FILE. Edit canonical lessons, then run npm run generate:books.
% canonical-source-hash: fnv1a64:b6274d53a6baa615
% canonical-lessons: GU-C525-silp, GU-C525-kotarvu, GU-C525-akar, GU-C525-madhvu, GU-C525-pradarsan

\chapter{Carving and Shaping}
\label{ch:gu-carving-and-shaping}
\hlchaptermodality{525}

\begin{chapteropening}
\textbf{By the end of this chapter:} \emph{I can talk about sculpture, carving and exhibitions.}

Complete the last lesson of chapter 525: I can talk about sculpture, carving and exhibitions.
\end{chapteropening}

\section[\texorpdfstring{\gu{શિલ્પ} (śilp)}{શિલ્પ (śilp)}]{\gu{શિલ્પ} (śilp) — a sculpture}

\label{lesson:GU-C525-silp}



\begin{quote}
Before the new one: say the Gujarati for hide-and-seek, then the Gujarati for a puzzle.
\end{quote}

\subsection*{You'll want to know: \gu{શિલ્પ}}
\textbf{\gu{શિલ્પ}} — \emph{śilp} — \textquotedblleft{}a sculpture\textquotedblright{}.

\begin{quote}
\gu{મંદિરમાં} \gu{પથ્થરનાં} \gu{ઘણાં} \gu{સુંદર} \gu{શિલ્પ} \gu{છે}.
\end{quote}

\begin{quote}
\emph{mandirmā̃ paththarnā̃ ghaṇā̃ sundar śilp che.}
\end{quote}

\begin{quote}
There are many beautiful stone sculptures in the temple.
\end{quote}

A neuter noun.

\begin{cousinweb}
\textbf{\gu{શિલ્પ}} (\emph{śilp}) is a tatsama, taken from Sanskrit \emph{śilpa}, \textquotedblleft{}an art, a craft\textquotedblright{} (Monier-Williams); in Gujarati it is a work of sculpture.
\end{cousinweb}

\begin{grammarlens}[title={What you've built}]
One more word for describing what you lived and felt.
\end{grammarlens}

\subsection*{Your turn}
\begin{itemize}
\raggedright
  \item \emph{Say it:} \emph{śilp}
  \item \emph{Say it:} \emph{śilp}, once more
  \item \emph{Say it:} \emph{mandirmā̃ paththarnā̃ ghaṇā̃ sundar śilp che.} — There are many beautiful stone sculptures in the temple.
  \item \emph{Recall:} say the Gujarati for hide-and-seek, then the Gujarati for a puzzle, then say \emph{śilp} again
\end{itemize}

\begin{tcolorbox}[breakable,colback=teal!4,colframe=teal!35!black,title={Before you move on}]
What does \gu{શિલ્પ} mean? (\textquotedblleft{}A sculpture\textquotedblright{}.) Say it once more.
\end{tcolorbox}
\section[\texorpdfstring{\gu{કોતરવું} (kotarvũ)}{કોતરવું (kotarvũ)}]{\gu{કોતરવું} (kotarvũ) — to carve}

\label{lesson:GU-C525-kotarvu}



\begin{quote}
Before the new one: say the Gujarati for a puzzle, then the Gujarati for a sculpture.
\end{quote}

\subsection*{You'll want to know: \gu{કોતરવું}}
\textbf{\gu{કોતરવું}} — \emph{kotarvũ} — \textquotedblleft{}to carve\textquotedblright{}.

\begin{quote}
\gu{કારીગરે} \gu{લાકડામાં} \gu{એક} \gu{મોર} \gu{કોતર્યો}.
\end{quote}

\begin{quote}
\emph{kārīgre lākḍāmā̃ ek mor kotaryo.}
\end{quote}

\begin{quote}
The craftsman carved a peacock in the wood.
\end{quote}

Past, with \gu{તેણે} (he or she …): \gu{તેણે} \gu{મોર} \gu{કોતર્યો}, \emph{teṇe mor kotaryo}.

\begin{cousinweb}
\textbf{\gu{કોતરવું}} (\emph{kotarvũ}) comes from \emph{*kōtr}, a form CDIAL reconstructs, and CDIAL lists the Gujarati verb, \textquotedblleft{}to dig, carve\textquotedblright{} (CDIAL 3512).
\end{cousinweb}

\begin{grammarlens}[title={What you've built}]
One more word for describing what you lived and felt.
\end{grammarlens}

\subsection*{Your turn}
\begin{itemize}
\raggedright
  \item \emph{Say it:} \emph{kotarvũ}
  \item \emph{Say it:} \emph{kotarvũ}, once more
  \item \emph{Say it:} \emph{kārīgre lākḍāmā̃ ek mor kotaryo.} — The craftsman carved a peacock in the wood.
  \item \emph{Recall:} say the Gujarati for a puzzle, then the Gujarati for a sculpture, then say \emph{kotarvũ} again
\end{itemize}

\begin{tcolorbox}[breakable,colback=teal!4,colframe=teal!35!black,title={Before you move on}]
What does \gu{કોતરવું} mean? (\textquotedblleft{}To carve\textquotedblright{}.) Say it once more.
\end{tcolorbox}
\section[\texorpdfstring{\gu{આકાર} (ākār)}{આકાર (ākār)}]{\gu{આકાર} (ākār) — a shape}

\label{lesson:GU-C525-akar}



\begin{quote}
Before the new one: say the Gujarati for a sculpture, then the Gujarati for to carve.
\end{quote}

\subsection*{You'll want to know: \gu{આકાર}}
\textbf{\gu{આકાર}} — \emph{ākār} — \textquotedblleft{}a shape\textquotedblright{}.

\begin{quote}
\gu{આ} \gu{પથ્થરનો} \gu{આકાર} \gu{ઘણો} \gu{સુંદર} \gu{છે}.
\end{quote}

\begin{quote}
\emph{ā paththarno ākār ghaṇo sundar che.}
\end{quote}

\begin{quote}
The shape of this stone is very beautiful.
\end{quote}

A masculine noun.

\begin{cousinweb}
\textbf{\gu{આકાર}} (\emph{ākār}) is a tatsama, taken from Sanskrit \emph{ākāra}, \textquotedblleft{}form, figure, shape\textquotedblright{} (Monier-Williams).
\end{cousinweb}

\begin{grammarlens}[title={What you've built}]
One more word for describing what you lived and felt.
\end{grammarlens}

\subsection*{Your turn}
\begin{itemize}
\raggedright
  \item \emph{Say it:} \emph{ākār}
  \item \emph{Say it:} \emph{ākār}, once more
  \item \emph{Say it:} \emph{ā paththarno ākār ghaṇo sundar che.} — The shape of this stone is very beautiful.
  \item \emph{Recall:} say the Gujarati for a sculpture, then the Gujarati for to carve, then say \emph{ākār} again
\end{itemize}

\begin{tcolorbox}[breakable,colback=teal!4,colframe=teal!35!black,title={Before you move on}]
What does \gu{આકાર} mean? (\textquotedblleft{}A shape\textquotedblright{}.) Say it once more.
\end{tcolorbox}
\section[\texorpdfstring{\gu{મઢવું} (maḍhvũ)}{મઢવું (maḍhvũ)}]{\gu{મઢવું} (maḍhvũ) — to frame (a picture)}

\label{lesson:GU-C525-madhvu}



\begin{quote}
Before the new one: say the Gujarati for to carve, then the Gujarati for a shape.
\end{quote}

\subsection*{You'll want to know: \gu{મઢવું}}
\textbf{\gu{મઢવું}} — \emph{maḍhvũ} — \textquotedblleft{}to frame (a picture)\textquotedblright{}.

\begin{quote}
\gu{અરુણે} \gu{મીરાનું} \gu{ચિત્ર} \gu{મઢ્યું}.
\end{quote}

\begin{quote}
\emph{Aruṇe Mīrānũ citra maḍhyũ.}
\end{quote}

\begin{quote}
Arun framed Mira's picture.
\end{quote}

Past, with \gu{તેણે} (he or she …): \gu{તેણે} \gu{ચિત્ર} \gu{મઢ્યું}, \emph{teṇe citra maḍhyũ}.

\begin{cousinweb}
\textbf{\gu{મઢવું}} (\emph{maḍhvũ}) comes from \emph{*maḍh}, a form CDIAL reconstructs, and CDIAL lists the Gujarati verb as \textquotedblleft{}to cover\textquotedblright{}, said of a book or a drum (CDIAL 9729); a picture is covered with glass and a frame.
\end{cousinweb}

\begin{grammarlens}[title={What you've built}]
One more word for describing what you lived and felt.
\end{grammarlens}

\subsection*{Your turn}
\begin{itemize}
\raggedright
  \item \emph{Say it:} \emph{maḍhvũ}
  \item \emph{Say it:} \emph{maḍhvũ}, once more
  \item \emph{Say it:} \emph{Aruṇe Mīrānũ citra maḍhyũ.} — Arun framed Mira's picture.
  \item \emph{Recall:} say the Gujarati for to carve, then the Gujarati for a shape, then say \emph{maḍhvũ} again
\end{itemize}

\begin{tcolorbox}[breakable,colback=teal!4,colframe=teal!35!black,title={Before you move on}]
What does \gu{મઢવું} mean? (\textquotedblleft{}To frame (a picture)\textquotedblright{}.) Say it once more.
\end{tcolorbox}
\section[\texorpdfstring{\gu{પ્રદર્શન} (pradarśan)}{પ્રદર્શન (pradarśan)}]{\gu{પ્રદર્શન} (pradarśan) — an exhibition}

\label{lesson:GU-C525-pradarsan}



\begin{quote}
Before the new one: say the Gujarati for a shape, then the Gujarati for to frame (a picture).
\end{quote}

\subsection*{You'll want to know: \gu{પ્રદર્શન}}
\textbf{\gu{પ્રદર્શન}} — \emph{pradarśan} — \textquotedblleft{}an exhibition\textquotedblright{}.

\begin{quote}
\gu{શાળામાં} \gu{ચિત્રોનું} \gu{પ્રદર્શન} \gu{છે}.
\end{quote}

\begin{quote}
\emph{śāḷāmā̃ citronũ pradarśan che.}
\end{quote}

\begin{quote}
There is an exhibition of paintings at the school.
\end{quote}

A neuter noun.

\begin{cousinweb}
\textbf{\gu{પ્રદર્શન}} (\emph{pradarśan}) is a tatsama, taken from Sanskrit \emph{pradarśana}, \textquotedblleft{}a look, a showing\textquotedblright{} (Monier-Williams).
\end{cousinweb}

\begin{grammarlens}[title={What you've built}]
One more word for describing what you lived and felt.
\end{grammarlens}

\subsection*{Your turn}
\begin{itemize}
\raggedright
  \item \emph{Say it:} \emph{pradarśan}
  \item \emph{Say it:} \emph{pradarśan}, once more
  \item \emph{Say it:} \emph{śāḷāmā̃ citronũ pradarśan che.} — There is an exhibition of paintings at the school.
  \item \emph{Recall:} say the Gujarati for a shape, then the Gujarati for to frame (a picture), then say \emph{pradarśan} again
\end{itemize}

\begin{tcolorbox}[breakable,colback=teal!4,colframe=teal!35!black,title={Before you move on}]
What does \gu{પ્રદર્શન} mean? (\textquotedblleft{}An exhibition\textquotedblright{}.) Say it once more.
\end{tcolorbox}
